\begin{frame}{자주 묻는 질문 (8.1)}
  \scriptsize
  \begin{tabular}{@{}p{0.3\textwidth}p{0.66\textwidth}@{}}
    \toprule
    질문 & 핵심 답 \\
    \midrule
    두 알고리즘 성능이 거의 같으면 & ``비슷하다''는 것 자체가 관찰 ---
      학습중/학습후를 \textbf{나눠} 보면 겉보기엔 비슷해도 다른
      트레이드오프라는 것을 발견할 수 있음 \\
    매번 결과가 달라 비교가 안 되면 & $\varepsilon$ 여러 개로 반복 실행한
      평균(과 가능하면 분산) 비교가 \textbf{표준적 방법} \\
    시드 몇 개가 정석인가 & 실행 비용 vs 결론 안정성의 절충 --- 표 기반
      환경은 시드당 몇 초$\sim$분 $\to$ \textbf{3$\sim$5개 현실적 하한},
      10개면 더 단단. 표준편차 1개 안에 겹치면 정직한 결론 =
      ``유의미한 차이 관측 불가'' \\
    코드를 처음부터 써야 하나 & 강의 코드 재용이 \textbf{전제} ---
      (a) 재사용한 부분(Ch\,6.3 갱신 규칙)과 (b) 팀이 수정·추가한
      부분(환경 인터페이스, 시드 처리, 평가 프로토콜)을 \textbf{구분} \\
    \midrule
    시간이 부족해 시드를 5$\to$3개로 줄였으면 & 성공 기준은 ``깨끗한
      숫자''가 아니라 \textbf{``정직한 절차''} --- ``시간 부족으로
      3시드로만 확인, SARSA 시드 간 편차가 커서 5시드 재확인 권장''을
      ``잘 안 된 부분''에 쓰는 것이 감추는 것보다 항상 낫다 \\
    \bottomrule
  \end{tabular}
\end{frame}
